% GENERATED FILE. Edit canonical lessons, then run npm run generate:books.
% canonical-source-hash: fnv1a64:2d6a0d7a03a0c7fe
% canonical-lessons: ML-C457-abhyartthikkuka, ML-C457-kramikarikkuka, ML-C457-nirvahikkuka, ML-C457-sthapikkuka, ML-C457-nilanirttuka

\chapter{Getting Things Done}
\label{ch:ml-getting-things-done}
\hlchaptermodality{457}

\begin{chapteropening}
\textbf{By the end of this chapter:} \emph{I can say to appeal, to arrange, to carry out, to found and to maintain.}

Complete the last lesson of chapter 457: I can say to appeal, to arrange, to carry out, to found and to maintain.
\end{chapteropening}

\section[\texorpdfstring{\ml{അഭ്യർത്ഥിക്കുക} (abhyartthikkuka)}{അഭ്യർത്ഥിക്കുക (abhyartthikkuka)}]{\ml{അഭ്യർത്ഥിക്കുക} (abhyartthikkuka) — to appeal; to request formally}

\label{lesson:ML-C457-abhyartthikkuka}



\begin{quote}
Before the new one: say the Malayalam for to sort, then the Malayalam for to divide.
\end{quote}

\subsection*{You'll want to know: \ml{അഭ്യർത്ഥിക്കുക}}
\textbf{\ml{അഭ്യർത്ഥിക്കുക}} — \emph{abhyartthikkuka} — \textquotedblleft{}to appeal; to request formally\textquotedblright{}.

\begin{quote}
\ml{വെള്ളം} \ml{പാഴാക്കരുതെന്ന്} \ml{പഞ്ചായത്ത്} \ml{ജനങ്ങളോട്} \ml{അഭ്യർത്ഥിച്ചു}.
\end{quote}

\begin{quote}
\emph{veḷḷaṁ pāḻākkarutennŭ pañcāyattŭ janaṅṅaḷōṭŭ abhyartthiccu.}
\end{quote}

\begin{quote}
The panchayat appealed to the people not to waste water.
\end{quote}

Past: \ml{അഭ്യർത്ഥിച്ചു}, \emph{abhyartthiccu}.

The person appealed to takes -\ml{ഓട്}. The noun is \ml{അഭ്യർത്ഥന}, an appeal.

\begin{cousinweb}
\textbf{\ml{അഭ്യർത്ഥിക്കുക}} (\emph{abhyartthikkuka}) is built from Sanskrit \emph{abhyarthana}, \textquotedblleft{}asking, requesting\textquotedblright{} (Monier-Williams), with the Malayalam verb ending -\ml{ിക്കുക}.
\end{cousinweb}

\begin{grammarlens}[title={What you've built}]
One more word for saying what happens if.
\end{grammarlens}

\subsection*{Your turn}
\begin{itemize}
\raggedright
  \item \emph{Say it:} \emph{abhyartthikkuka}
  \item \emph{Say it:} \emph{abhyartthikkuka}, once more
  \item \emph{Say it:} \emph{veḷḷaṁ pāḻākkarutennŭ pañcāyattŭ janaṅṅaḷōṭŭ abhyartthiccu.} — The panchayat appealed to the people not to waste water.
  \item \emph{Recall:} say the Malayalam for to sort, then the Malayalam for to divide, then say \emph{abhyartthikkuka} again
\end{itemize}

\begin{tcolorbox}[breakable,colback=teal!4,colframe=teal!35!black,title={Before you move on}]
What does \ml{അഭ്യർത്ഥിക്കുക} mean? (\textquotedblleft{}To appeal; to request formally\textquotedblright{}.) Say it once more.
\end{tcolorbox}
\section[\texorpdfstring{\ml{ക്രമീകരിക്കുക} (kramīkarikkuka)}{ക്രമീകരിക്കുക (kramīkarikkuka)}]{\ml{ക്രമീകരിക്കുക} (kramīkarikkuka) — to set in order; to arrange}

\label{lesson:ML-C457-kramikarikkuka}



\begin{quote}
Before the new one: say the Malayalam for to divide, then the Malayalam for to appeal.
\end{quote}

\subsection*{You'll want to know: \ml{ക്രമീകരിക്കുക}}
\textbf{\ml{ക്രമീകരിക്കുക}} — \emph{kramīkarikkuka} — \textquotedblleft{}to set in order; to arrange\textquotedblright{}.

\begin{quote}
\ml{അവൾ} \ml{മുറിയിലെ} \ml{കസേരകൾ} \ml{ക്രമീകരിച്ചു}.
\end{quote}

\begin{quote}
\emph{avaḷ muṟiyile kasērakaḷ kramīkariccu.}
\end{quote}

\begin{quote}
She arranged the chairs in the room.
\end{quote}

Past: \ml{ക്രമീകരിച്ചു}, \emph{kramīkariccu}.

The noun is \ml{ക്രമീകരണം}, an arrangement.

\begin{cousinweb}
\textbf{\ml{ക്രമീകരിക്കുക}} (\emph{kramīkarikkuka}) is built from Sanskrit \emph{krama}, \textquotedblleft{}a step, a series, an order\textquotedblright{} (CDIAL 3577), with -\ml{ീകരിക്കുക}, \textquotedblleft{}to make into\textquotedblright{}: to put in order.
\end{cousinweb}

\begin{grammarlens}[title={What you've built}]
One more word for saying what happens if.
\end{grammarlens}

\subsection*{Your turn}
\begin{itemize}
\raggedright
  \item \emph{Say it:} \emph{kramīkarikkuka}
  \item \emph{Say it:} \emph{kramīkarikkuka}, once more
  \item \emph{Say it:} \emph{avaḷ muṟiyile kasērakaḷ kramīkariccu.} — She arranged the chairs in the room.
  \item \emph{Recall:} say the Malayalam for to divide, then the Malayalam for to appeal, then say \emph{kramīkarikkuka} again
\end{itemize}

\begin{tcolorbox}[breakable,colback=teal!4,colframe=teal!35!black,title={Before you move on}]
What does \ml{ക്രമീകരിക്കുക} mean? (\textquotedblleft{}To set in order; to arrange\textquotedblright{}.) Say it once more.
\end{tcolorbox}
\section[\texorpdfstring{\ml{നിർവഹിക്കുക} (nirvahikkuka)}{നിർവഹിക്കുക (nirvahikkuka)}]{\ml{നിർവഹിക്കുക} (nirvahikkuka) — to carry out; to perform}

\label{lesson:ML-C457-nirvahikkuka}



\begin{quote}
Before the new one: say the Malayalam for to appeal, then the Malayalam for to set in order.
\end{quote}

\subsection*{You'll want to know: \ml{നിർവഹിക്കുക}}
\textbf{\ml{നിർവഹിക്കുക}} — \emph{nirvahikkuka} — \textquotedblleft{}to carry out; to perform\textquotedblright{}.

\begin{quote}
\ml{അവൻ} \ml{തന്റെ} \ml{ജോലി} \ml{നന്നായി} \ml{നിർവഹിച്ചു}.
\end{quote}

\begin{quote}
\emph{avan tanṟe jōli nannāyi nirvahiccu.}
\end{quote}

\begin{quote}
He carried out his work well.
\end{quote}

Past: \ml{നിർവഹിച്ചു}, \emph{nirvahiccu}.

Said of a duty, a task or a ceremony.

\begin{cousinweb}
\textbf{\ml{നിർവഹിക്കുക}} (\emph{nirvahikkuka}) is built from Sanskrit \emph{nirvahaṇa}, \textquotedblleft{}the end, completion\textquotedblright{} (Monier-Williams), with the Malayalam verb ending -\ml{ിക്കുക}: to carry a thing to its end.
\end{cousinweb}

\begin{grammarlens}[title={What you've built}]
One more word for saying what happens if.
\end{grammarlens}

\subsection*{Your turn}
\begin{itemize}
\raggedright
  \item \emph{Say it:} \emph{nirvahikkuka}
  \item \emph{Say it:} \emph{nirvahikkuka}, once more
  \item \emph{Say it:} \emph{avan tanṟe jōli nannāyi nirvahiccu.} — He carried out his work well.
  \item \emph{Recall:} say the Malayalam for to appeal, then the Malayalam for to set in order, then say \emph{nirvahikkuka} again
\end{itemize}

\begin{tcolorbox}[breakable,colback=teal!4,colframe=teal!35!black,title={Before you move on}]
What does \ml{നിർവഹിക്കുക} mean? (\textquotedblleft{}To carry out; to perform\textquotedblright{}.) Say it once more.
\end{tcolorbox}
\section[\texorpdfstring{\ml{സ്ഥാപിക്കുക} (sthāpikkuka)}{സ്ഥാപിക്കുക (sthāpikkuka)}]{\ml{സ്ഥാപിക്കുക} (sthāpikkuka) — to establish; to found}

\label{lesson:ML-C457-sthapikkuka}



\begin{quote}
Before the new one: say the Malayalam for to set in order, then the Malayalam for to carry out.
\end{quote}

\subsection*{You'll want to know: \ml{സ്ഥാപിക്കുക}}
\textbf{\ml{സ്ഥാപിക്കുക}} — \emph{sthāpikkuka} — \textquotedblleft{}to establish; to found\textquotedblright{}.

\begin{quote}
\ml{അവന്റെ} \ml{മുത്തച്ഛൻ} \ml{ആ} \ml{കട} \ml{സ്ഥാപിച്ചു}.
\end{quote}

\begin{quote}
\emph{avanṟe muttacchan ā kaṭa sthāpiccu.}
\end{quote}

\begin{quote}
His grandfather founded that shop.
\end{quote}

Past: \ml{സ്ഥാപിച്ചു}, \emph{sthāpiccu}.

Also to set up a thing in a place: \ml{പ്രതിമ} \ml{സ്ഥാപിച്ചു}, they set up a statue.

\begin{cousinweb}
\textbf{\ml{സ്ഥാപിക്കുക}} (\emph{sthāpikkuka}) is built from Sanskrit \emph{sthāpana}, \textquotedblleft{}causing to stand\textquotedblright{} (CDIAL 13755), with the Malayalam verb ending -\ml{ിക്കുക}: to make a thing stand.
\end{cousinweb}

\begin{grammarlens}[title={What you've built}]
One more word for saying what happens if.
\end{grammarlens}

\subsection*{Your turn}
\begin{itemize}
\raggedright
  \item \emph{Say it:} \emph{sthāpikkuka}
  \item \emph{Say it:} \emph{sthāpikkuka}, once more
  \item \emph{Say it:} \emph{avanṟe muttacchan ā kaṭa sthāpiccu.} — His grandfather founded that shop.
  \item \emph{Recall:} say the Malayalam for to set in order, then the Malayalam for to carry out, then say \emph{sthāpikkuka} again
\end{itemize}

\begin{tcolorbox}[breakable,colback=teal!4,colframe=teal!35!black,title={Before you move on}]
What does \ml{സ്ഥാപിക്കുക} mean? (\textquotedblleft{}To establish; to found\textquotedblright{}.) Say it once more.
\end{tcolorbox}
\section[\texorpdfstring{\ml{നിലനിർത്തുക} (nilanirttuka)}{നിലനിർത്തുക (nilanirttuka)}]{\ml{നിലനിർത്തുക} (nilanirttuka) — to maintain; to retain}

\label{lesson:ML-C457-nilanirttuka}



\begin{quote}
Before the new one: say the Malayalam for to carry out, then the Malayalam for to establish.
\end{quote}

\subsection*{You'll want to know: \ml{നിലനിർത്തുക}}
\textbf{\ml{നിലനിർത്തുക}} — \emph{nilanirttuka} — \textquotedblleft{}to maintain; to retain\textquotedblright{}.

\begin{quote}
\ml{അവൾ} \ml{ഒന്നാം} \ml{സ്ഥാനം} \ml{നിലനിർത്തി}.
\end{quote}

\begin{quote}
\emph{avaḷ onnāṁ sthānaṁ nilanirtti.}
\end{quote}

\begin{quote}
She retained first place.
\end{quote}

Past: \ml{നിലനിർത്തി}, \emph{nilanirtti}.

The intransitive is \ml{നിലനിൽക്കുക}, to last.

\begin{cousinweb}
\textbf{\ml{നിലനിർത്തുക}} (\emph{nilanirttuka}) is a native Malayalam compound: \ml{നില}, \textquotedblleft{}a standing, a position\textquotedblright{}, and \ml{നിർത്തുക}, \textquotedblleft{}to make stand, to stop\textquotedblright{}, which DEDR writes \emph{niṟuttuka}; both belong to the root \emph{nil}, \textquotedblleft{}to stand\textquotedblright{} (DEDR 3675): to keep a thing standing.
\end{cousinweb}

\begin{grammarlens}[title={What you've built}]
One more word for saying what happens if.
\end{grammarlens}

\subsection*{Your turn}
\begin{itemize}
\raggedright
  \item \emph{Say it:} \emph{nilanirttuka}
  \item \emph{Say it:} \emph{nilanirttuka}, once more
  \item \emph{Say it:} \emph{avaḷ onnāṁ sthānaṁ nilanirtti.} — She retained first place.
  \item \emph{Recall:} say the Malayalam for to carry out, then the Malayalam for to establish, then say \emph{nilanirttuka} again
\end{itemize}

\begin{tcolorbox}[breakable,colback=teal!4,colframe=teal!35!black,title={Before you move on}]
What does \ml{നിലനിർത്തുക} mean? (\textquotedblleft{}To maintain; to retain\textquotedblright{}.) Say it once more.
\end{tcolorbox}
